\chapter{El departamento que no está en los expedientes}
\label{cap:ausencias}

Este libro tiene veintiocho capítulos sobre un lugar donde viven doce mil trescientas personas,
y las personas casi no aparecen.

Aparecen como firmantes de resoluciones, como partes de un remate, como titulares de una
matrícula, como damnificados de un aluvión que otro relevamiento entrevistó. No aparecen
comprando, estudiando, enfermándose, yéndose ni quedándose. \textbf{El apéndice \ref{cap:fuentes} lo advierte: el departamento aparece tal como el Estado lo escribió. Este capítulo hace el
inventario de lo que esa elección dejó afuera.}

No es un capítulo de conclusiones. Es un mapa de huecos, con lo poco que el archivo sí dejó de
cada uno, y con lo que haría falta para llenarlos. Sirve para que un lector sepa qué no está
leyendo, y para que quien continúe este trabajo no empiece de cero.

\section{Por qué faltan}

La razón no es una omisión del autor: es una propiedad del archivo.

\textbf{Los expedientes registran la propiedad y el agua porque de eso se ocupan.} Un remate
deja constancia de quién compró y en cuánto, no de quién trabajaba esa tierra. Una determinación
de línea de ribera fija coordenadas, no cuántas casas quedan adentro. Un decreto de promoción
turística consigna un beneficiario, no cuántos puestos de trabajo generó ---y el capítulo
\ref{cap:trabajo} muestra que nadie está obligado a informarlo---.

\textbf{De modo que la ausencia también es un dato.} Un territorio del que el Estado guarda
tres siglos de títulos y casi ningún registro de vida cotidiana está diciendo algo sobre qué
consideró digno de registro.

\begin{cautela}
Con una salvedad que hay que hacer, porque si no este capítulo se vuelve una queja.

\textbf{Los registros de la vida cotidiana existen y no están en los archivos que este libro
consultó.} Hay legajos escolares, historias clínicas, libros parroquiales, padrones de
comercios, expedientes de la comisaría. No están en el Boletín Oficial ni en la Dirección
General de Inmuebles, y varios ---con razón--- no son públicos.

Lo que falta, entonces, no es que el Estado no registre la vida. Es que \textbf{este
relevamiento eligió un archivo y ese archivo no la contiene}.
\end{cautela}

\section{Lo que sí quedó, por dimensión}

\subsection*{Los nombres que se borraron}

Hay una ausencia que este capítulo puede medir, y es nueva. El capítulo \ref{cap:siglo} reúne las
doce nóminas con que el Estado provincial nombró las unidades de este departamento entre 1909 y
1940 ---policía, Código Rural, circuitos electorales--- y las cruza con los dos nomencladores
vigentes: el campo «finca» del catastro parcelario y la Base de Asentamientos Humanos.

\textbf{Veinticuatro de los cincuenta y dos nombres no están en ninguno de los dos} (lámina
\ref{fig:parajesnom}). Siete de ellos son lugares donde el circuito electoral de 1918 declara que
vive gente que vota ---Despensilla, Corral de Barrancas, El Churcal, El Monte, Las Garrapatas, El
Túnel y La Despensa---, y otros los registraban todavía los circuitos de 1921 y el padrón de 1940: El
Bordo, la Estación del ferrocarril, Severino, Los Chusos, Higueritas, Los Matos, Nogalito y La
Sierra.

\begin{cautela}
\textbf{No se afirma que esos lugares hayan desaparecido.} Un topónimo puede seguir vivo en el uso
y no estar en un nomenclador, y este libro no fue al territorio a preguntarlo ---que es,
exactamente, la carencia que este capítulo inventaría---. Lo que queda establecido es más chico y
se puede verificar: \textbf{el Estado dejó de registrarlos, y no hay acto que lo disponga}. Se
borraron entre dos formularios.
\end{cautela}

\subsection*{La escuela}

Es la dimensión con más huellas, y son huellas indirectas: las escuelas aparecen porque tienen
inmuebles.

El capítulo \ref{cap:tierrafiscal} documenta la \textbf{Escuela Nº 4118 «Dr.\ Gustavo Martínez
Zuviría»}, del paraje El Gallinato, que funciona desde 1968 en un terreno ajeno y cuyo suelo la
Provincia adquirió por prescripción en 2015. \textbf{Y hay dos nombres entre los que este libro transcribe del archivo temprano, y los dos aparecen por jubilarse, no por enseñar}: en julio de 1911 el Ministerio de Hacienda jubila a \textbf{Jacoba Luquez de Gómez, directora de la escuela elemental de la parroquia} (capítulo \ref{cap:presencia}), y el decreto 1.653 del 20 de abril de 1918 jubila a \textbf{Juana Luque, Directora de la Escuela de la Caldera}, con ciento cuatro pesos cincuenta centavos mensuales (capítulo \ref{cap:siglo}). La semejanza de los apellidos no autoriza a identificarlas entre sí ni a suponer parentesco: son dos actos distintos, con siete años de distancia. Una versión anterior de este párrafo decía que el nombre de 1918 era el único; el propio libro transcribía el de 1911. El recuento vale sobre lo que el libro transcribe y no sobre el archivo: los tramos leídos después de los primeros no se repasaron buscando nombres de docentes. Aparecen además la \textbf{Escuela Nº 136 Juana
Moro de López}, la \textbf{Escuela Nº 660 de Finca Mojotoro}, la \textbf{Escuela Nº 4534 de Los Yacones} y el \textbf{Colegio Nº 5045 «Senado Provincial»}.

\textbf{Lo que no se sabe:} cuántos alumnos, cuántos establecimientos hay en total, qué
proporción termina la primaria y dónde cursan la secundaria quienes la terminan. Y algo que el
capítulo \ref{cap:poblacion} vuelve pertinente: si los chicos del corredor van a la escuela del
departamento o viajan a la capital.

\subsection*{La salud}

Es el hueco más grande del libro. Lo que hay es histórico y escaso: un lazareto improvisado
durante la viruela de 1908 (capítulo \ref{cap:siglo}), una estación sanitaria prometida en 1935
que no dejó rastro y un guarda sanitario a sueldo provincial (capítulo \ref{cap:presencia}), y
una deuda del municipio con Sanidad en 1940 (capítulo \ref{cap:hacienda}). \textbf{Del presente,
nada.}

No se sabe cuántos centros de salud hay, con qué horarios, con qué personal, ni a dónde va
alguien que necesita atención de urgencia a las tres de la mañana en un departamento a veinte
kilómetros de la capital. Tampoco si el crecimiento poblacional que el capítulo
\ref{cap:poblacion} documenta tuvo correlato en infraestructura sanitaria.

\subsection*{El comercio y el abastecimiento}

Siete menciones, todas de paso.

No se sabe dónde se compra la comida, cuántos comercios hay, si existe feria, ni qué proporción
del consumo se hace en la capital. En un departamento cuya superficie agropecuaria declarada cayó un ochenta y ocho por ciento ---capítulo \ref{cap:tierra}---, saber de dónde viene hoy la comida
no es un detalle costumbrista.

\subsection*{El transporte}

Acá el archivo dejó algo mejor, y por una vía inesperada.

El capítulo \ref{cap:resistencias} documenta que \textbf{Transporte Hnos.\ Mendoza S.R.L.}
tramitó ante el Juzgado de Minas la renovación de la concesión de la cantera de áridos
«Macarena». \textbf{La misma razón social aparece como transportista y como concesionaria de
extracción sobre el cauce.} El capítulo \ref{cap:redes} documenta además la creación de una
autoridad metropolitana de transporte, que el capítulo \ref{cap:aguabaja} señala como la única
figura de coordinación que la provincia construyó para este corredor.

\textbf{Lo que no se sabe:} qué servicios de pasajeros existen hoy, con qué frecuencia y a qué
costo, y si alcanzan a los parajes o solo al casco urbano.

\subsection*{Los jóvenes y la migración}

El capítulo \ref{cap:resistencias} registra, tomándolo del libro de la
comunidad, el \textbf{despoblamiento del cerro y la migración de los jóvenes}. Y el capítulo
\ref{cap:tierra} documenta el movimiento inverso y anterior: la llegada de trabajadores del
tabaco desde el interior provincial y desde Bolivia, primero golondrinas y después radicados.

\textbf{Dos migraciones en direcciones opuestas, separadas por décadas, y de ninguna se sabe
cuánta gente movió.}

\subsection*{La vida asociativa}

Es la dimensión mejor conservada, porque las asociaciones sacan personería y las personerías se
publican.

El capítulo \ref{cap:resistencias} reúne al menos siete asociaciones civiles con personería en el departamento; de varias, este libro no pudo establecer la continuidad después de 2015.

\textbf{Y acá el hueco se llenó mientras este capítulo se escribía}, porque las asambleas de
entidades civiles se publican en el Boletín Oficial y forman una serie.

\begin{ficha}{Entidades del departamento con asamblea publicada, 2016--2026}
\textbf{Clubes:} Social y Deportivo de Caza y Pesca La Caldera; Deportivo Jardín; Social y
Deportivo San Roque; Deportivo La Verbena; Social y Deportivo La Caldera; Social, Cultural y
Deportivo Sportivo Vaqueros; Club Nuevo Horizonte.

\textbf{Jubilados y cultura:} Centro de Jubilados y Pensionados Cristo Monumental, con
asambleas en 2017, 2018, 2019, 2021, 2022, 2023, 2024 y 2025; Biblioteca Popular Roberto
Romero, de Vaqueros, con serie igual de continua; Centro Gaucho El Palenque.

\textbf{Emergencias:} Asociación de Bomberos Voluntarios 2 de Junio, Vaqueros; y Asociación
Civil Papa Francisco, Bomberos Voluntarios de La Caldera, con asamblea en 2026 ---y, en marzo de ese año, intervenida: la Inspección General de Personas Jurídicas encontró su sede cerrada y abandonada, declaró la caducidad de los mandatos de su comisión directiva y nombró interventoras (Resolución 26/26, B.O.\ Nº 22.162; capítulo \ref{cap:amparo}). Fundada en septiembre de 2013 según el registro del Sistema Nacional de Bomberos Voluntarios, con personería de 2016; su cuartel está en la avenida Costanera, en su cruce con la avenida San Martín, en el pueblo de La Caldera---.

\textbf{Producción y trabajo:} Asociación de Pequeños Productores Agropecuarios; Cooperativa de
Trabajo Salta Norte 3 Ltda.; \textbf{Cámara de Turismo del Departamento La Caldera}.

\textbf{Vecinales y urbanizaciones:} Centro Vecinal Barrio El Durazno; Club de Campo Vertientes
Eco-Pueblo, de Vaqueros; Asociación para el Progreso de los Argentinos.
\end{ficha}

\textbf{No es un departamento sin vida asociativa: es un departamento cuya vida asociativa este
libro no había buscado.} Y estaba en la misma publicación que usó durante todo el
relevamiento, bajo un tipo de aviso que no había mirado.

\begin{cautela}
\textbf{Qué prueba esta serie y qué no.} Prueba que esas entidades existen, que tienen
personería y que convocaron asamblea en las fechas que se indican. \textbf{No prueba cuántos
socios tienen, ni qué hacen, ni si la asamblea llegó a celebrarse.}

Y es una lección sobre el propio método, que este capítulo haría mal en callar: \textbf{la
ausencia que el libro atribuía al archivo era, en parte, una ausencia de búsqueda.} El Boletín
Oficial registra bastante más vida cotidiana de la que este relevamiento supuso, sólo que en
secciones que no consultó.

Cuánto más hay en esas secciones ---concursos, sucesorios, edictos judiciales--- este libro no
lo sabe. \textbf{El resto de este capítulo debe leerse con esa advertencia puesta: no todo lo
que declara faltante falta.}
\end{cautela}

\subsection*{La iglesia}

Aparece por sus bienes y sus obras, no por su vida: la ley de 1935 que destinó siete mil pesos
a reformar la Iglesia Parroquial (capítulos \ref{cap:defensas} y \ref{cap:presencia}), la
posesión treintañal que la Curia promovió en 1939 (capítulo \ref{cap:fincas}), el loteo de una
congregación en Vaqueros (capítulo \ref{cap:vaqueros}) y la pavimentación, en el Plan de Obras
Públicas de 1974, de la ruta de acceso al \textbf{Cristo Redentor}.

No hay en este libro una sola línea sobre la vida parroquial de un departamento donde el
monumento religioso es el principal atractivo turístico. \textbf{Los libros parroquiales
---bautismos, matrimonios, defunciones--- son la serie demográfica más antigua y continua que
existe en cualquier pueblo argentino}, y este relevamiento no los consultó.

\subsection*{Las mujeres}

\textbf{Casi ninguna mención como tales.}

No es que el libro las haya excluido: es que el archivo que consultó las registra como
titulares de dominio o por el cargo que ocupan, y poco más. María del Valle Esteban aparece como propietaria de la matrícula 3.161; María Isabel Gómez, en un expediente de 1926; Silvia Santamaría firma edictos del organismo hídrico desde 2009; una senadora, una jueza, una abogada de Fiscalía, una cacique y una intendenta interina aparecen por su función. El único dato que las mira como población está en el diagnóstico social del Plan de Manejo: el 58 por ciento de los jefes de hogar encuestados son mujeres (capítulo \ref{cap:amparo}).

\textbf{De la vida de las mujeres de este departamento ---qué trabajan, cómo se organizan, qué
pasa con las que quedan cuando los jóvenes se van--- este libro no sabe absolutamente nada}, y
no puede saberlo por la vía que eligió.

\textbf{Lo único que puede aportar es demografía}, y conviene ponerla porque acota lo que puede
suponerse.

\begin{ficha}{Censo 2022, población por sexo y edad --- La Caldera}
\textbf{6.125 mujeres y 6.174 varones}, que suman el total definitivo de 12.299: un índice de masculinidad de \textbf{100,8 varones cada cien mujeres}. Una primera versión de este cuadro, tomada de un procesamiento parcial, daba 6.092 y 6.145; el cuadro por sexo del INDEC la corrige. La población está equilibrada.

Por grupos, sobre las 12.246 personas en viviendas particulares: 0 a 4 años, 6,5\,\%; 5 a 14, 16,2\,\%; 15 a 19, 7,3\,\%; \textbf{20 a 29, 15,2\,\%};
30 a 39, 15,2\,\%; 40 a 49, 15,5\,\%; 50 a 59, 10,5\,\%; 60 a 69, 7,9\,\%; 70 y más, 5,8\,\%.

\textbf{La población de 15 a 64 años es el 68,1\,\%} del total, y la de 0 a 14, el 22,6.

\textit{Procesamiento Redatam por radio censal, edades simples, septiembre de 2026. Una versión anterior de este cuadro, tomada de un procesamiento parcial, daba 5,0 por ciento de 0 a 4 años y 68,9 de 15 a 64.}
\end{ficha}

\textbf{Dos cosas se siguen de ahí, y la segunda pone un límite a un testimonio que este libro
recoge.}

\textbf{No hay desequilibrio de sexos.} En departamentos con emigración masculina de temporada,
o con actividades extractivas que atraen varones, el índice se despega de cien. Acá no.

\textbf{Y en el departamento no hay hueco en la franja joven, pero en una de sus fracciones
sí.} El capítulo \ref{cap:resistencias} registra, tomado del libro de la
comunidad, el despoblamiento del cerro y la migración de los jóvenes. En la pirámide
departamental esa migración no se ve: el grupo de 20 a 29 años pesa casi lo mismo que el de 30 a
39 y el de 40 a 49. \textbf{Desagregada por fracción, aparece.}

\begin{ficha}{Censo 2022, estructura de edades por fracción censal}
\textbf{Fracción 66077-01} ---7.735 personas---: 0 a 14, 22,1\,\%; \textbf{15 a 29, 22,4\,\%};
30 a 44, 22,6\,\%; 45 a 64, 23,7\,\%; 65 y más, 9,2\,\%.

\textbf{Fracción 66077-02} ---3.747---: 0 a 14, 24,2\,\%; \textbf{15 a 29, 24,5\,\%}; 30 a 44,
23,5\,\%; 45 a 64, 18,9\,\%; 65 y más, 8,9\,\%.

\textbf{Fracción 66077-03} ---764---: 0 a 14, 20,4\,\%; \textbf{15 a 29, 14,4\,\%}; 30 a 44,
26,8\,\%; 45 a 64, 26,2\,\%; \textbf{65 y más, 12,2\,\%}.
\end{ficha}

\textbf{La fracción 66077-03 tiene ocho a diez puntos menos de población joven que las otras
dos}, y a la vez la mayor proporción de adultos de 30 a 64 y de mayores de 65. Es la misma
fracción que tiene el 43,17 por ciento de sus viviendas sin residentes permanentes y la peor
cobertura de agua de red del departamento ---58 por ciento de hogares sin conexión, según el
capítulo \ref{cap:aguabaja}---.

\begin{cautela}
\textbf{Una primera lectura de este capítulo acercó este dato al testimonio}, y la distancia importa.

\textbf{Lo que se ve} es una fracción con menos jóvenes, más adultos mayores, más viviendas
vacías y menos servicios que el resto del departamento. Parecía compatible con el despoblamiento del cerro que el libro de la comunidad describe. La cartografía censal, que sigue, obliga a separar las dos cosas.
\end{cautela}

\textbf{Y ahora sí se puede afirmar qué territorio es cada fracción.} El capítulo
\ref{cap:poblacion} lo suponía por magnitud demográfica; la cartografía censal lo confirma.

\begin{ficha}{Fracciones censales del departamento, cartografía del organismo censal}
\textbf{Fracción 66077-01}, sector occidental y meridional: \textbf{es el municipio de
Vaqueros}.

\textbf{Fracción 66077-02}, sector central: \textbf{contiene el pueblo de La Caldera y La
Calderilla}.

\textbf{Fracción 66077-03}, sector oriental: \textbf{contiene El Gallinato y la Escuela Nº
4118}.

\textbf{Dos comprobaciones independientes coinciden.} La geométrica ---los puntos
georreferenciados del pueblo, de La Calderilla, de El Gallinato y de la escuela caen dentro de
los polígonos indicados--- y la aritmética: la fracción 01 registra 7.735 personas en viviendas
particulares y el municipio de Vaqueros, \textbf{7.735} en viviendas particulares; las fracciones 02 y 03 suman
\textbf{4.511}, que es exactamente la población en viviendas particulares del municipio de La Caldera. \textit{(Las poblaciones totales son 7.774 y 4.525; la diferencia son las viviendas colectivas, capítulo \ref{cap:poblacion}.)}
\end{ficha}

\begin{cautela}
\textbf{De modo que la fracción con menos jóvenes, más viviendas vacías y peor cobertura de agua
es el sector oriental del departamento: el de El Gallinato.} No es el cerro que el relevamiento
comunitario menciona ---que está al oeste, sobre el faldeo--- sino el otro extremo.

\textbf{Eso desarma la coincidencia que este apartado insinuaba, y conviene decirlo sin
rodeos.} El despoblamiento que el capítulo \ref{cap:resistencias} recoge se refiere al cerro,
que cae en la fracción 01 o en la 02. \textbf{El hueco demográfico está en la 03, que es El
Gallinato.} Son dos fenómenos distintos y este libro los estaba superponiendo.

Lo que la fracción 03 muestra ---pocos jóvenes, muchas viviendas sin residentes permanentes, la
peor cobertura de red--- describe un sector rural de baja densidad con vivienda de fin de
semana, que es precisamente donde el apéndice \ref{cap:pedidos} ubica los dos fraccionamientos de El Gallinato que el capítulo \ref{cap:opacidad} observa en el visor, todavía sin identificar. No deben confundirse con las parcelas en cinta del capítulo \ref{cap:loteo}, que están al sur del pueblo, en la fracción 02.

\textbf{Y sigue sin saberse qué pasó con quienes se fueron del cerro}, porque el cerro está en
otra fracción y su despoblamiento no se ve en estos números.

\textbf{Y son 764 personas.} En un universo de ese tamaño, un puñado de casos mueve varios
puntos porcentuales. Las proporciones de la 03 hay que leerlas con esa reserva.
\end{cautela}

\begin{cautela}
\textbf{Nada de esto refuta el testimonio, y conviene decir por qué.} El libro de la comunidad habla del
\textbf{cerro}, no del departamento. Un éxodo desde la zona alta hacia el casco urbano o hacia
los loteos del valle no aparecería en una pirámide departamental: la gente se mueve adentro del
mismo polígono censal.

\textbf{Lo que la pirámide sí descarta} es que el departamento en su conjunto esté perdiendo
jóvenes. Con 58,4 por ciento de crecimiento poblacional en doce años ---capítulo
\ref{cap:poblacion}--- y una estructura de edades sin vacíos, \textbf{La Caldera recibe
población, no la expulsa}.

La pregunta que queda, y que este capítulo no puede contestar, es \textbf{quién se fue del cerro
y adónde}.
\end{cautela}

\section{Qué haría falta}

Ninguna de estas dimensiones se llena con más expedientes. Todas se llenan con las mismas tres
cosas, y en este orden.

\textbf{Primero, las series que sí existen y no se pidieron.} Padrón de establecimientos
educativos y matrícula por año; nómina de efectores de salud y su dotación; padrón de comercios
habilitados del municipio; nómina de servicios de transporte autorizados; y los libros
parroquiales, que son los únicos que cubren dos siglos. Están en organismos que atienden al
público y la mayoría no es reservada.

\textbf{Segundo, el procesamiento del censo.} Población por rama de actividad, condición de
ocupación, nivel educativo alcanzado y lugar de trabajo. Es un pedido al organismo censal, no
una investigación.

\textbf{Y tercero, lo único que no tiene sustituto: preguntar.} Un protocolo de entrevistas ---preparado aparte y no incluido en este volumen--- sirve para las dimensiones que el libro ya documenta. Para éstas haría
falta otro, con otras personas: docentes, personal de salud, comerciantes, transportistas, y
sobre todo mujeres de las tres generaciones que el corredor tiene.

\begin{cautela}
\textbf{Una advertencia sobre el orden, porque se puede hacer mal.}

Las tres cosas anteriores se pueden hacer en cualquier secuencia salvo una: \textbf{preguntar
antes de haber pedido}. Si alguien va a entrevistar a una docente sin haber pedido el padrón de
matrícula, va a terminar usando el recuerdo de esa persona como si fuera una estadística, que es
exactamente lo que este libro se prohíbe en el capítulo \ref{cap:metodo}.

El testimonio sirve para lo que la serie no puede decir: por qué, cómo se vivió, qué se decidió.
No para cuántos.
\end{cautela}

\section{Por qué este capítulo está acá y no es un prólogo de otro libro}

Podría haberse omitido. Un libro no está obligado a inventariar lo que no hizo, y ningún lector
le reclamaría a un trabajo sobre el uso del suelo que además explicara el sistema de salud.

Está por dos razones.

\textbf{La primera es de método.} Este libro sostiene, capítulo tras capítulo, que las ausencias
documentales son informativas: que un anexo no publicado, una coordenada remitida a la sede o un
padrón que no existe dicen algo. \textbf{Sería incoherente aplicar ese criterio a los
organismos y no a sí mismo.}

\textbf{Y la segunda es práctica.} Lo que falta acá es, casi exactamente, el contenido de un
segundo volumen: el mismo departamento contado por quienes viven en él en lugar de por quienes
lo escrituraron. Este capítulo es su esqueleto, y queda publicado a propósito, \textbf{para que
la lista de lo que falta sea pública antes de que exista el libro que la llene}.

